\section{本章小结}
\label{sec:budget-summary}

本章把采样生成、误差分析、方差控制与梯度表示放到同一本账上，回答“给定任务和预算，怎样安排一次随机计算”。回到章首的问题：计算$I=\mathbb{P}(X>4)$，要求$95\%$把握下相对误差不超过$10\%$。$I\approx0.0183$离零较远，相对要求有意义，按式~\eqref{eq:budget-mixed-tolerance}给一个绝对下限以防目标接近零即可。直接平均约需$20\,600$个样本；先用几十个试运行样本拟合$\operatorname{Exp}(\lambda)$建议分布，再做重要性采样，约需$1990$个正式样本，成本约$3000$。两者都在$10^{5}$的预算之内，后者约省$85\%$；正式样本数按两阶段方法由试运行确定，最终估计只用正式样本。

\paragraph{误差与成本的口径}
误差要求先要说清是均方误差、指定概率下的绝对误差，还是相对误差；目标可能为零时用混合容差。无偏方法之间比较工作量归一化方差$\sigma^{2}c$，它在固定预算与固定精度两种提法下给出同一排序；有偏方法还要选择控制偏差的参数，近似链的例子中最优安排使均方根误差只按$B^{-1/3}$下降。成本要计入被拒绝的候选、矩阵分解与试运行等一次性开销，以及链的初始化；相关链按有效样本计成本，而有效样本数依赖被积函数，抽稀只节省存储。

\paragraph{预算的分配}
试运行的最优规模按$\sqrt{C}$增长，占预算的比例随预算增大而缩小，最终误差对它的具体取值通常不敏感。分层的最优分配在考虑成本时为$n_k\propto p_k\sigma_k/\sqrt{c_k}$。对单峰且混合良好的链，一条长链的均方误差最小，多条短链各自带着初值偏差；多次启动的价值在于暴露链没有走遍目标。成组随机化中，组内点数$n$决定精度，组数$R$决定误差估计的可靠性，不能把全部点当作独立样本估计误差。

\paragraph{方法的选择}
可用的方法由可计算的信息决定：分位数、条件分布、有限包络、密度比、期望已知的辅助函数、可微的生成映射或得分。在可用的方法之间，以同一任务的总成本比较：一维双峰例子中直接生成比最好的Metropolis便宜十几倍；方差控制的组合要重新计算，对偶变量之后再用控制变量$U$毫无收益，而换成对称的控制函数则收益巨大；梯度估计没有普遍的优劣，被积函数振荡时得分函数方差小，维数高时重参数化方差小，被积函数不连续时重参数化可能静默地给出错误结果。

\paragraph{试运行与停止的条件}
试运行只用于决定，正式计算用新的随机数，固定方法的结论才在条件意义下适用。MCMC在试运行中调参、正式计算中固定转移核；按轨迹持续换核，即使每个核都保持目标不变，整体也可能收敛到错误的分布。稀有贡献、权重集中与链间分歧都能让有限记录显得平稳，经验检查只能发现问题，不能证明正确。依据误差估计决定何时停止会使偏斜被积函数的区间系统性地失效，只在要求足够严时渐近可靠；把定样本数与计算估计分开的两阶段方法保留了固定样本数的结论。

最终报告一次随机计算，应分别给出计算偏差（及其来源，如近似生成或初值）、随机波动（及其估计方式和所依据的条件）以及实际成本。这里的“精度”是对一个确定的数做随机计算时的精度：分布已知、积分待算，随机性完全来自我们所用的随机数，可以通过增加计算来减小。它与由有限数据推断未知总体时的精度是两回事，后者的误差来自数据本身，加大计算量无法消除，相应的理论见\BookChapterRef{chap:parametric-statistical-inference}{卷八第1章“参数化统计推断”}。
